%% kapa: exam=1 round=27 year=2016 period=1 form=A subject=real_estate subject_ko=부동산학원론 seq=2 number=82
%% source: Q-Net 감정평가사 기출문제 (cst00309) artlSeq=5169402 file=extracted/1008/5169402/제27회 감정평가사 1차 문제/제27회 감정평가사 1교시문제(A형).pdf
%% conversion: kordoc 4.24.51 -> tools/gen1.py (auto); review=proofread-sonnet; flags=none
\begin{question}{82}
용도지역지구제에 관한 설명으로 옳지 않은 것은?
\begin{choices}
  \item 토지이용에 수반되는 부(-)의 외부효과를 제거하거나 감소시키는 것에 목적이 있다.
  \item 국토의 계획 및 이용에 관한 법령상 시ㆍ도지사는 도시의 자연환경 및 경관을 보호하고 도시민에게 건전한 여가ㆍ휴식공간을 제공하기 위하여 도시지역 안에서 식생이 양호한 산지의 개발을 제한할 필요가 있다고 인정하면 도시자연공원구역의 지정을 도시ㆍ군관리계획으로 결정할 수 있다.
  \item 사적시장이 외부효과에 대한 효율적인 해결책을 제시하지 못할 경우, 정부가 부동산규제의 한 방법으로 채택할 수 있다.
  \item 국토의 계획 및 이용에 관한 법령상 시ㆍ도지사는 시설보호지구를 문화재, 중요 시설물 및 문화적ㆍ생태적으로 보존가치가 큰 지역의 보호와 보존에 필요하면 그 용도지구의 지정 또는 변경을 도시ㆍ군관리계획으로 결정한다.
  \item 국토의 계획 및 이용에 관한 법령상 시ㆍ도지사는 도시ㆍ군관리계획결정으로 주거지역을 세분하여 지정할 수 있는데, 공동주택 중심의 양호한 주거환경을 보호하기 위하여 필요한 지역은 제2종전용주거지역으로 지정할 수 있다.
\end{choices}
\end{question}
\answer{4}
